% Chapter 02: Authoring Schemas
% Defining custom step types as YAML overlays on core schemas

\chapter{Authoring Schemas}
\label{ch:authoring-schemas}

% ============================================================
\section{What Kit Schemas Are}
\label{sec:what-schemas-are}
% ============================================================

A Kit schema is a JSON Schema (Draft 2020-12) document that defines
domain-specific step types as \emph{overlays} on top of the gert core schema.
The core schema defines the base step types (\texttt{cli}, \texttt{tool},
\texttt{manual}, \texttt{branch}, \texttt{iterate}) and their required fields.
A Kit schema adds new step types with domain-specific field structures.

Key principles:

\begin{itemize}
  \item A Kit schema is \textbf{additive}. It does not modify or remove core
    step types. It adds new types that will be lowered into core types.
  \item A Kit schema is \textbf{validated at parse time}. When a runbook
    declares \texttt{apiVersion: runbook/v1+compliance/v1}, the gert parser
    loads the Kit schema and validates the runbook against it.
  \item A Kit schema is \textbf{domain-specific}. It should feel natural for
    authors working in that domain. A compliance Kit might define
    \texttt{approval-request} and \texttt{evidence-capture}; a workflow Kit
    might define \texttt{select-priority} and \texttt{escalate}.
\end{itemize}

% ============================================================
\section{Defining a Custom Step Type}
\label{sec:custom-step-type}
% ============================================================

A custom step type is defined as a JSON Schema object that extends the base
step schema with additional required and optional fields.

\paragraph{Schema overlay pattern:}

\begin{minted}{json}
{
  "$schema": "https://json-schema.org/draft/2020-12/schema",
  "$id": "https://gert.dev/schemas/kits/compliance-v1.json",
  "title": "Compliance Kit v1 Schema",
  "description": "Compliance-focused authoring vocabulary for gert",

  "definitions": {
    "approval-request": {
      "type": "object",
      "required": ["id", "type", "approvers"],
      "properties": {
        "id": { "type": "string" },
        "type": { "const": "approval-request" },
        "description": { "type": "string" },
        "approvers": {
          "type": "array",
          "items": { "type": "string" },
          "minItems": 1,
          "description": "List of user identifiers who can approve this step"
        },
        "timeout": {
          "type": "string",
          "pattern": "^[0-9]+(s|m|h|d)$",
          "description": "Max wait time (e.g., '2h', '30m'). Auto-reject on timeout."
        },
        "evidence-required": {
          "type": "boolean",
          "default": false,
          "description": "If true, approver must attach evidence"
        },
        "x-kit:metadata": {
          "type": "object",
          "description": "Kit-specific metadata preserved through lowering"
        }
      },
      "additionalProperties": false
    }
  },

  "oneOf": [
    { "$ref": "#/definitions/approval-request" }
  ]
}
\end{minted}

\paragraph{Field conventions:}

\begin{description}
  \item[\texttt{id}] --- REQUIRED. Unique step identifier within the runbook.
  \item[\texttt{type}] --- REQUIRED. Step type discriminator. Use \texttt{"const"}
    to enforce exact match.
  \item[\texttt{description}] --- Optional but recommended. Human-readable step
    summary.
  \item[Domain fields] --- Add required and optional fields specific to your
    domain (e.g., \texttt{approvers}, \texttt{timeout}, \texttt{evidence-required}).
  \item[\texttt{x-kit:*}] --- Prefix for Kit-defined metadata fields. Use this
    namespace for fields that should be preserved through lowering but are not
    part of the core schema.
\end{description}

% ============================================================
\section{Field Extension Conventions}
\label{sec:field-conventions}
% ============================================================

\paragraph{Required vs. optional fields:}

Use \texttt{"required"} to declare fields that MUST be present for semantic
correctness. Example: \texttt{approvers} is required for
\texttt{approval-request} because an approval without approvers is meaningless.

Optional fields should have sensible defaults or be truly optional. Document
the default behavior in the schema \texttt{description}.

\paragraph{Type constraints:}

Use JSON Schema \texttt{type}, \texttt{pattern}, \texttt{minItems},
\texttt{enum}, etc., to enforce structural correctness. Example:

\begin{minted}{json}
"timeout": {
  "type": "string",
  "pattern": "^[0-9]+(s|m|h|d)$",
  "description": "Duration with unit suffix (s=seconds, m=minutes, h=hours, d=days)"
}
\end{minted}

\paragraph{The \texttt{x-kit:} annotation prefix:}

Fields prefixed with \texttt{x-kit:} are Kit-specific metadata that will be
preserved through lowering but are not interpreted by the gert core. Use this
for:

\begin{itemize}
  \item Hints to the lowering compiler (e.g., \texttt{x-kit:compile-mode: strict})
  \item Metadata for projections (e.g., \texttt{x-kit:audit-category: sensitive})
  \item Domain-specific annotations (e.g., \texttt{x-kit:compliance-framework: SOC2})
\end{itemize}

The gert core ignores \texttt{x-kit:*} fields. Your Kit compiler can read them
from the AST and use them to inform lowering behavior.

% ============================================================
\section{Validation at Parse Time vs. Plan Time}
\label{sec:parse-vs-plan}
% ============================================================

There are two phases of validation in the gert pipeline:

\begin{description}
  \item[Parse-time validation] runs when \texttt{gert validate} is invoked or
    when a runbook is loaded. The parser validates the runbook YAML against
    the Kit schema (structural validation). This catches syntax errors, missing
    required fields, type mismatches, etc.

  \item[Plan-time validation] runs after lowering, when the planner is about to
    produce an ExecutionPlan. This is where semantic validators
    (\S\ref{ch:validation-rules}) run. Plan-time validators can see the full
    lowered runbook and enforce cross-step invariants (e.g., ``cleanup must
    follow provision'').
\end{description}

\paragraph{What belongs in the schema (parse-time):}

\begin{itemize}
  \item Field presence (required vs. optional)
  \item Field types (string, number, array, object)
  \item Format constraints (patterns, enums, ranges)
  \item Per-step structural correctness
\end{itemize}

\paragraph{What belongs in semantic validators (plan-time):}

\begin{itemize}
  \item Cross-step invariants (``every X must be followed by Y'')
  \item Business rules (``timeout must not exceed SLA'')
  \item Context-dependent validation (``if field A is set, field B is required'')
\end{itemize}

Use JSON Schema for structure; use semantic validators for semantics.

% ============================================================
\section{Example: \texttt{evidence-capture} Step Type}
\label{sec:example-evidence}
% ============================================================

Let's define a second custom step type: \texttt{evidence-capture}. This step
pauses execution and prompts a human to provide structured evidence (text,
file upload, or both).

\begin{minted}{json}
{
  "definitions": {
    "evidence-capture": {
      "type": "object",
      "required": ["id", "type", "prompt"],
      "properties": {
        "id": { "type": "string" },
        "type": { "const": "evidence-capture" },
        "description": { "type": "string" },
        "prompt": {
          "type": "string",
          "description": "Instruction text shown to the operator"
        },
        "evidence-type": {
          "type": "string",
          "enum": ["text", "attachment", "both"],
          "default": "text",
          "description": "What kind of evidence to collect"
        },
        "required-attestation": {
          "type": "boolean",
          "default": true,
          "description": "If true, operator must attest evidence is accurate"
        },
        "x-kit:audit-tag": {
          "type": "string",
          "description": "Tag for grouping evidence in audit reports"
        }
      },
      "additionalProperties": false
    }
  }
}
\end{minted}

\paragraph{Authoring example:}

A runbook author using the compliance Kit might write:

\begin{minted}{yaml}
apiVersion: runbook/v1+compliance/v1
meta:
  title: "Quarterly access review"
  kit:
    name: gert.compliance
    version: "^1.0.0"

steps:
  - id: request-approval
    type: approval-request
    description: "Get approval from security team"
    approvers:
      - alice@acme.example
      - bob@acme.example
    timeout: 2h
    evidence-required: true

  - id: capture-review-results
    type: evidence-capture
    description: "Upload access review spreadsheet"
    prompt: "Attach the completed access review CSV file"
    evidence-type: attachment
    x-kit:audit-tag: "access-review-q1-2025"
\end{minted}

This Kit-authored runbook will be validated against the compliance Kit schema
at parse time, then lowered into core step types before execution.

% ============================================================
\section{Schema Composition and Reuse}
\label{sec:schema-composition}
% ============================================================

For complex Kits with many step types, use JSON Schema \texttt{\$ref} and
\texttt{definitions} to avoid repetition:

\begin{minted}{json}
{
  "definitions": {
    "timeout-field": {
      "type": "string",
      "pattern": "^[0-9]+(s|m|h|d)$"
    },
    "user-list": {
      "type": "array",
      "items": { "type": "string" },
      "minItems": 1
    },
    "approval-request": {
      "type": "object",
      "required": ["id", "type", "approvers"],
      "properties": {
        "id": { "type": "string" },
        "type": { "const": "approval-request" },
        "approvers": { "$ref": "#/definitions/user-list" },
        "timeout": { "$ref": "#/definitions/timeout-field" }
      }
    },
    "escalate": {
      "type": "object",
      "required": ["id", "type", "escalate-to"],
      "properties": {
        "id": { "type": "string" },
        "type": { "const": "escalate" },
        "escalate-to": { "$ref": "#/definitions/user-list" },
        "timeout": { "$ref": "#/definitions/timeout-field" }
      }
    }
  }
}
\end{minted}

This reduces duplication and ensures consistency across step types.
